% XSEC: Just-in-Time Rule Injection design note
% Compile: pdflatex rule-injection.tex && pdflatex rule-injection.tex

\documentclass[11pt]{article}

\usepackage[margin=1in]{geometry}
\usepackage[utf8]{inputenc}
\usepackage[T1]{fontenc}
\usepackage{hyperref}
\usepackage{url}
\usepackage{booktabs}
\usepackage{amsmath}
\usepackage{amsfonts}
\usepackage{microtype}
\usepackage{enumitem}
\usepackage{graphicx}
\usepackage{xcolor}

\hypersetup{
  colorlinks=true,
  linkcolor=blue!60!black,
  citecolor=blue!60!black,
  urlcolor=blue!60!black,
}

\title{Just-in-Time Rule Injection for Autonomous Security Agents}

\author{
  XSEC Labs \\
  \texttt{github.com/uncesaii/xsec}
}

\date{}

\begin{document}

\maketitle

\begin{abstract}
Autonomous security agents accumulate guidance through several channels: a
static system prompt, model-pulled skills, and dynamically injected
methodology documents. All of these operate at the granularity of a
\emph{topic} --- a vulnerability class, a full technique, a target profile ---
and all of them are assembled either before generation begins or at a
coarse reflection checkpoint. None of them delivers a small, atomic
\emph{do this / do not do that} correction at the precise moment the agent is
about to take a specific wrong action. We survey the design space of
agent rule systems, including Oh My Pi's Time-Traveling Stream Rules (TTSR),
which is unique in triggering on a regular-expression or structural match
against the \emph{live output stream} and injecting mid-generation. We then
document the design decision for a rule-injection layer in XSEC that sits on
top of its existing dynamic playbooks and just-in-time (JIT) skills. The XSEC
layer keeps the key TTSR insight --- match the action, not just the aftermath
--- but adapts it to a turn-based native-API loop: a pure, LLM-free
\texttt{selectRules} function fires on the \emph{tool input} (the command or
diff about to run), scoped by phase, language, file glob, and tool, and injects
a one-line rule at the post-dispatch reflection checkpoint. We give the rule
schema, the selection logic, the injection and telemetry plumbing, five seed
security rules, and the rationale for why this belongs beside --- not inside ---
the playbook and skill layers.
\end{abstract}

\section{Introduction and Motivation}

A capable security agent already knows a great deal in the abstract. What it
lacks, repeatedly and expensively, is the right small reminder at the right
small moment: rotate your headers before you retry against a WAF; keep a
\texttt{Session} so the authenticated cookie survives the next request; do not
dump the whole credentials table when a single row proves the finding. Each of
these is one sentence. Each is worthless five turns early and five turns late.
And each is tied not to a topic the agent is thinking about, but to a concrete
\emph{action} the agent is about to emit --- a shell command, a code diff, an
HTTP request.

XSEC today injects guidance through two dynamic mechanisms and one static one,
described in Section~\ref{sec:gap}. All three deliver large documents keyed to a
topic. This note argues for, and specifies, a fourth mechanism at the opposite
end of the granularity axis: atomic, action-triggered rules. The contribution
is a design decision, not a benchmark. We (i) locate the idea in the prior art,
(ii) identify the specific gap in XSEC's current stack, and (iii) fix the schema,
selection function, injection point, telemetry, and controls precisely enough to
implement without further design work.

\section{Background: Oh My Pi's Time-Traveling Stream Rules}
\label{sec:ttsr}

Oh My Pi (OMP) introduces \emph{Time-Traveling Stream Rules}
(TTSR)~\cite{omp,ompttsr,ompwiki}, the closest prior art to what we propose and
the source of its central idea. A TTSR rule is an object with the fields:
\texttt{name}; \texttt{condition}, a regular expression evaluated over the output
stream; \texttt{astCondition}, an ast-grep structural pattern for matching code;
\texttt{globs}, a path scope; \texttt{interruptMode} and \texttt{repeatMode},
which govern firing behavior; and \texttt{text}, the content injected on a match.

The distinguishing property is \emph{when} matching runs. Rules are evaluated on
every streaming \texttt{message\_update} inside the agent session, not at request
assembly. Three matchers cooperate: \texttt{checkDelta()} runs the regex over the
buffered text and thinking as it streams; \texttt{checkSnapshot()} and
\texttt{checkAstSnapshot()} run the regex and ast-grep patterns over the
reconstructed source of an in-progress edit or write. When a rule matches and its
mode allows interruption, the agent \emph{aborts mid-token}, waits roughly
50\,ms, and retries generation from the same point --- now with the rule injected
as a
\texttt{<system-interrupt reason="rule\_violation" rule="\{name\}" path="\{path\}">}
block (a \texttt{<system-reminder ...>} variant is used for non-interrupting
matches against tool source). That abort-and-retry-from-the-same-position is the
``time travel'': the model re-generates the offending span having already been
told what is wrong with it, before the bad token ever reaches a tool.

Two further properties matter for our design. Injections survive context
compaction, because each is persisted as a hidden \texttt{ttsr-injection} custom
message rather than living only in the transient stream. And the whole subsystem
is gated: \texttt{ttsr.enabled} toggles it, \texttt{ttsr.disabledRules} and
\texttt{ttsr.builtinRules} select the active set, and any rule can opt out of
interruption with \texttt{interruptMode: never}. The user-facing signal is a
single amber ``Injecting rule: \guillemotright{}\,name'' banner in the TUI. OMP
keeps this stream-triggered layer distinct from \emph{static} context rules,
which are loaded into the system prompt from \texttt{.omp/rules/},
\texttt{AGENTS.md}, \texttt{.clinerules}, and \texttt{.cursorrules}.

\section{Prior Art in Agent Rule Systems}
\label{sec:priorart}

Most coding-agent rule systems share an architecture: rules are authored as
files with some metadata, selected by a coarse trigger, and merged into the
prompt at \emph{request-assembly} time --- before the model generates. Cursor's
\texttt{.cursor/rules} uses \texttt{.mdc} files whose frontmatter carries
\texttt{description}, \texttt{globs}, and \texttt{alwaysApply}, yielding four
rule types: Always, Auto-Attached (by glob), Agent-Requested (the model decides
from the description), and Manual~\cite{cursorrules}. Claude Code loads
\texttt{CLAUDE.md} plus \texttt{@import} references. OpenCode reads
\texttt{instructions} and \texttt{AGENTS.md}. Continue's \texttt{.continue/rules}
selects by glob.

\begin{table}[ht]
\centering
\small
\begin{tabular}{llll}
\toprule
System & Trigger & Injection time & Granularity \\
\midrule
Cursor \texttt{.cursor/rules} & glob / description / always & request assembly & file/topic \\
Claude Code \texttt{CLAUDE.md} & always (+ \texttt{@import}) & request assembly & project \\
OpenCode \texttt{AGENTS.md} & always & request assembly & project \\
Continue \texttt{.continue/rules} & glob & request assembly & file \\
OMP TTSR & regex / AST on live stream & mid-generation & span/action \\
\textbf{XSEC rules (proposed)} & regex on tool input & post-dispatch turn & action \\
\bottomrule
\end{tabular}
\caption{Trigger, injection time, and granularity across agent rule systems.
Everyone except OMP triggers on globs, file-in-context, or a model decision from
a description, and injects before generation. OMP is unique in triggering on a
match against the live output stream and injecting mid-generation. XSEC's
proposed layer keeps the action-level trigger but moves it to the tool input at a
turn boundary.}
\end{table}

The common axis is clear from the table. Cursor, Claude Code, OpenCode, and
Continue all trigger on \emph{globs}, on \emph{file-in-context}, or on a
\emph{model decision from a description}, and all inject at request-assembly
time. OMP alone triggers on a live-stream match and injects mid-generation.

Underneath every one of these systems is the same economy that motivates the
tool-search tool~\cite{toolsearch}: context is scarce and expensive, so
guidance should be disclosed progressively rather than loaded all at once. A
single always-on ruleset bloats the prompt, dilutes attention, and pays for
rules that never apply. The rule mechanisms above are, in effect, progressive
disclosure of rules --- surface a rule only when its trigger says it is relevant.
XSEC's design takes the same principle to the finest grain: surface one line only
at the moment its specific action is imminent.

\section{XSEC Today, and the Gap}
\label{sec:gap}

XSEC has two JIT-ish dynamic layers. Both trigger identically: a
\texttt{matchCount >= 2} threshold of regex hits over roughly the twenty most
recent tool-result strings.

\paragraph{Dynamic playbooks.} In
\texttt{packages/core/src/agent/playbooks.ts}, a \texttt{PLAYBOOKS} map holds
per-vulnerability-class methodology. \texttt{detectPlaybooks()} scores the recent
tool results, \texttt{buildPlaybookInjection()} renders the winning playbook, and
\texttt{native-loop.ts} injects it at the 30\% turn-budget checkpoint. A playbook
is roughly 3.6k tokens, emits \texttt{playbook\_injected} telemetry, and is gated
by the \texttt{dynamicPlaybooks} feature flag.

\paragraph{JIT skills.} In \texttt{packages/core/src/agent/skills/}, each
\texttt{SkillDefinition} carries \texttt{\{id, name, description, triggers}
(a regex array)\texttt{, estimated\_tokens, content\}}. Skills are surfaced to the
model through \texttt{list\_skills}/\texttt{load\_skill} and pulled on demand by
the model itself, gated by the \texttt{jitSkills} feature flag.

Both mechanisms deliver \emph{large methodology documents}. Playbooks are keyed
to a vulnerability class; skills are full methods the model chooses to read. The
static system prompt is larger still. What none of the three provides is a
\emph{small atomic do/don't rule delivered at the moment of the offending
action}, scoped by file glob, language, specific tool, and phase, and --- the
decisive point --- triggered on the tool \emph{input} (the command or diff about
to execute) rather than on tool results after the fact. That is the gap.

\section{The XSEC Design}
\label{sec:design}

We adopt the TTSR insight --- correct the action, not the aftermath --- but
adapt it to XSEC's execution model. XSEC's agent loop is turn-based over a native
API; it does not stream into an interruptible token buffer the way OMP does.
A mid-token abort-and-retry is therefore a v2 concern (Section~\ref{sec:future}).
The v1 design triggers on tool input and injects at the existing turn boundary.

\subsection{Rule format}

Rules live in \texttt{packages/core/src/agent/rules/*.yaml}. Each rule is atomic:

\begin{itemize}[leftmargin=*,itemsep=1pt]
  \item \texttt{id} --- stable identifier, used in telemetry and \texttt{disabledRules}.
  \item \texttt{name} --- short human label for the TUI banner.
  \item \texttt{severity} --- \texttt{hint} (silent inject) or \texttt{warn} (inject plus amber TUI line).
  \item \texttt{trigger} --- the narrowing conjunction: \texttt{phase: [recon|exploit|report]}, \texttt{lang: [...]}, \texttt{tool: [...]}, \texttt{glob: [...]}, \texttt{regex: [...]}.
  \item \texttt{rule} --- a one-line DO/DON'T string, the injected payload.
  \item \texttt{repeat} --- \texttt{once-per-scan}, \texttt{on-change}, or \texttt{always}.
\end{itemize}

\subsection{Selection logic}

Selection is a \emph{pure function}, \texttt{selectRules(context)}, with no LLM
call. Its context is
\texttt{\{phase, editedPath, toolName, toolInput, recentToolResultTexts\}}. It
filters the rule set as a cascade of cheap conjunctions:

\begin{enumerate}[leftmargin=*,itemsep=1pt]
  \item \textbf{phase} --- drop rules whose \texttt{phase} does not include the current phase.
  \item \textbf{lang / glob} --- match \texttt{glob} against \texttt{editedPath}; drop on mismatch.
  \item \textbf{tool} --- drop rules whose \texttt{tool} does not include \texttt{toolName}.
  \item \textbf{regex} --- run \texttt{trigger.regex} over \texttt{toolInput} --- the command or diff about to run.
\end{enumerate}

The key insight is in the last step: matching the \emph{input} means the rule
fires on the action the agent is about to take, not merely on what came back.
Unlike the playbook/skill \texttt{matchCount >= 2} heuristic over results, a
\emph{single} regex hit is sufficient here, because the preceding phase, glob, and
tool filters have already narrowed the candidate set to rules that are relevant
by construction. A lone match at that point is signal, not noise.

\subsection{Injection point}

Rules inject at the \emph{post-tool-dispatch, pre-next-turn} point in
\texttt{native-loop.ts} --- the same reflection checkpoint that already hosts
playbook injection. This is deliberately \emph{not} a mid-token abort: XSEC's
loop is turn-based, so the natural seam is the turn boundary. The payload matches
the codebase's existing convention, reusing the reminder block shape rather than
inventing a new one:

\begin{quote}\ttfamily
<system-reminder reason="rule" rule="\{id\}" path="\{path\}">\{rule\}</system-reminder>
\end{quote}

Compaction survival reuses the mechanism playbooks already rely on: a
\texttt{messagesContainText(marker)} check re-injects the rule if compaction has
dropped it from the working context, so a fired rule does not silently vanish
across a summarization boundary.

\subsection{Telemetry and controls}

Every firing emits \texttt{onEvent("rule\_injected", \{turn, ruleId, phase,
trigger, reason\})} and a corresponding \texttt{logEvent} to the database,
exactly paralleling the existing \texttt{playbook\_injected} path. For
\texttt{warn}-severity rules the TUI prints the amber ``Injecting rule:
\guillemotright{}\,name'' line, mirroring OMP's UX; \texttt{hint} rules inject
silently.

Controls follow the established feature-flag pattern. A master flag,
\texttt{features.ruleInjection}, defaults \textbf{OFF}. \texttt{rules.disabledRules}
suppresses named rules and \texttt{rules.builtinRules} selects the active builtin
set, matching the \texttt{ttsr.disabledRules}/\texttt{ttsr.builtinRules} shape from
OMP and the flag conventions already used by \texttt{dynamicPlaybooks} and
\texttt{jitSkills}.

\subsection{Why a fourth layer}

The three layers already present and the one proposed occupy distinct points on
the granularity axis, and that is precisely why the new layer earns its place
rather than folding into an existing one:

\begin{itemize}[leftmargin=*,itemsep=1pt]
  \item \textbf{Playbooks} --- vulnerability-class methodology; target-level; large document; triggered on results.
  \item \textbf{Skills} --- a full method; model-pulled on demand; large document.
  \item \textbf{Rules} --- atomic ``you are about to do this specific thing wrong, right now''; action-triggered on tool input; roughly one line.
\end{itemize}

A rule is not a small playbook, and a playbook is not a verbose rule. They
trigger on different signals (input vs.\ results), at different grains (a line
vs.\ a document), for different purposes (correct an action vs.\ teach a class).
Collapsing them would either bloat every action-level correction into a document
or dilute methodology into fragments.

\section{Seed Rules}
\label{sec:seed}

We seed the builtin set with five rules drawn from recurring, expensive mistakes
observed in security-agent runs. Each is one line; each is scoped so it fires
only when its specific action is imminent.

\begin{enumerate}[leftmargin=*,itemsep=2pt]
  \item \textbf{py-exploit-session} --- when a Python exploit issues sequential
  authenticated requests, use a \texttt{requests.Session} so cookie/auth state
  carries across calls instead of being dropped each request.
  \item \textbf{waf-rotate-headers} --- on a 403 or WAF block, rotate the
  User-Agent and headers before retrying rather than hammering the same
  fingerprint.
  \item \textbf{shell-injection-safe-poc} --- build proof-of-concept subprocess
  calls with an argv list, not \texttt{shell=True}, so the PoC does not itself
  introduce shell injection.
  \item \textbf{minimal-proof-no-exfil} --- extract the minimal, non-PII proof
  that demonstrates the finding; do not dump entire credential or user tables.
  \item \textbf{idor-body-and-path} --- when testing IDOR, change the identifier
  in \emph{both} the URL path and the request body, and confirm the response is
  the same resource, i.e.\ that an access boundary was actually crossed.
\end{enumerate}

\section{Decision}
\label{sec:decision}

We will add a rule-injection layer to XSEC as specified in
Section~\ref{sec:design}: YAML rule files under
\texttt{packages/core/src/agent/rules/}, a pure LLM-free \texttt{selectRules}
function triggering on tool input through a phase\,$\rightarrow$\,lang/glob\,$\rightarrow$\,tool\,$\rightarrow$\,regex
cascade, injection at the post-dispatch reflection checkpoint in
\texttt{native-loop.ts} using the existing \texttt{<system-reminder>} convention,
compaction survival via the \texttt{messagesContainText} re-inject trick,
\texttt{rule\_injected} telemetry paralleling \texttt{playbook\_injected}, and the
five seed rules of Section~\ref{sec:seed}. The layer ships behind
\texttt{features.ruleInjection}, defaulting OFF, and is complementary to --- not a
replacement for --- dynamic playbooks and JIT skills. We deliberately adopt
TTSR's action-level trigger while declining its mid-token abort mechanism, whose
cost is not justified by XSEC's turn-based loop at v1.

\section{Limitations and Future Work}
\label{sec:future}

The v1 design has two structural limits, both inherited from the decision to
trigger at the turn boundary rather than mid-generation. First, a rule fires
\emph{after} the offending tool has been dispatched: the correction lands before
the \emph{next} turn, so a genuinely destructive first action is not prevented,
only flagged and corrected forward. A v2 pre-execution interception hook --- run
\texttt{selectRules} against a proposed tool call before dispatch, and block or
rewrite on a match --- would close this window and bring XSEC closer to TTSR's
prevent-before-emit guarantee. Second, v1 matches only regexes over the tool
input string; it cannot reason about code structure. Adding \emph{AST triggers}
in the spirit of TTSR's \texttt{astCondition}/\texttt{checkAstSnapshot()} would
let rules match structural patterns in a proposed diff (e.g.\ a
\texttt{subprocess} call with \texttt{shell=True} regardless of formatting)
rather than brittle text patterns. Both are deferred to keep v1 a pure,
side-effect-free selection function that is cheap to test and safe to ship OFF by
default.

\begin{thebibliography}{9}

\bibitem{omp}
Oh My Pi.
\newblock \url{https://github.com/can1357/oh-my-pi}.

\bibitem{ompttsr}
Oh My Pi X --- TTSR Injection Lifecycle.
\newblock \url{https://github.com/rohreid/oh-my-pi-x/blob/main/docs/ttsr-injection-lifecycle.md}.

\bibitem{ompwiki}
Oh My Pi --- System Prompts (DeepWiki).
\newblock \url{https://deepwiki.com/can1357/oh-my-pi/4.4-system-prompts}.

\bibitem{cursorrules}
Cursor --- Rules for AI.
\newblock \url{https://docs.cursor.com/context/rules}.

\bibitem{toolsearch}
Anthropic --- Tool Search Tool.
\newblock \url{https://platform.claude.com/docs/en/agents-and-tools/tool-use/tool-search-tool}.

\end{thebibliography}

\end{document}
